\chapter{Outputs and Execution}\label{r01:ch:outputs-execution}

\code{create_kernel} returns a \cls{Kernel} object. This chapter explains what that
object contains, how to read the generated code, how the just-in-time compiler turns it
into a Python function, and how other tools embed it into C++ programs.

\section{The \texttt{Kernel} object}\label{r01:sec:kernel-object}

\cls{Kernel} and its subclass \cls{GpuKernel} are defined in
\file{codegen/kernel.py}. They are plain containers produced by the last stage of the
code generator:
\begin{center}\small
\begin{tabularx}{\linewidth}{@{}l>{\raggedright\arraybackslash}X@{}}
  \toprule
  \textbf{Member} & \textbf{Content}\\
  \midrule
  \code{name} & function name (settable; default \code{"kernel"})\\
  \code{target} & the resolved \code{Target}\\
  \code{parameters} & tuple of \cls{Parameter} objects, sorted by name\\
  \code{body} & the kernel body as a backend AST (\cref{r01:ch:ir})\\
  \code{required_headers} & set of headers to include (\code{<cmath>}, \code{"pystencils_runtime/x86.hpp"}, \dots)\\
  \code{metadata} & dictionary, e.g.\ \texttt{\{'ghost\_layers': 1\}}\\
  \code{get_fields()} & the set of \code{Field}s accessed\\
  \code{get_c_code()} & the C/C++ function as text (without includes)\\
  \code{get_ir_code()} & the same, printed as typed IR\\
  \code{compile()} & JIT-compile and return a callable \code{KernelWrapper}\\
  \code{get_launch_configuration()} & (\code{GpuKernel} only) a fresh launch configuration\\
  \bottomrule
\end{tabularx}
\end{center}
For the Jacobi kernel of \cref{r01:ch:getting-started}:
\outsnippet[caption={\texttt{str(kernel)}}]{ex01_first_kernel.str}
\outsnippet[caption={\texttt{kernel.required\_headers}}]{ex01_first_kernel.headers}
\outsnippet[caption={\texttt{kernel.metadata}}]{ex01_first_kernel.metadata}

\subsection{Parameters and their properties}

Every kernel parameter is a \cls{Parameter} (\file{codegen/parameters.py}) with a
\code{name}, a \code{dtype} and a set of \term{properties} that say what the parameter
\emph{means}:
\outsnippet[nolol=false,caption={Name, type and properties of every parameter of the Jacobi kernel.},label=r01:lst:params]{ex01_first_kernel.params}
The properties (\file{codegen/properties.py}) are frozen dataclasses:
\begin{description}
  \item[\code{FieldBasePtr(field)}] the parameter is the data pointer of \code{field};
  \item[\code{FieldShape(field, coordinate)}] it is the size of \code{field} in one
        coordinate;
  \item[\code{FieldStride(field, coordinate)}] it is a stride, in \emph{elements} (not
        bytes);
  \item[(none)] a scalar parameter from a free symbol (such as \code{alpha} in
        \cref{r01:lst:heat-code}), or a reduction target pointer.
\end{description}
This metadata is what makes kernels usable outside \ps. A runtime system only has to
look at the properties to know which array, size or stride to pass for each argument.
The CPU JIT (\cref{r01:sec:jit}) and pystencils-sfg (\cref{r01:sec:external-use}) both work this
way. Convenience accessors are \code{param.fields}, \code{param.is_field_parameter} and
\code{param.get_properties(FieldShape)}.

\section{Anatomy of the generated code}\label{r01:sec:anatomy}

Every kernel is a single C function \code{void <name>(<parameters>)}. CUDA and HIP
kernels carry the \code{__global__} prefix. The names follow fixed conventions from
\file{src/pystencils/defaults.py}:
\begin{center}\small
\begin{tabular}{@{}ll@{}}
  \toprule
  \textbf{Item} & \textbf{Name}\\
  \midrule
  data pointer of field \code{f} & \code{_data_f}\\
  size and stride of \code{f} in coordinate $i$ & \code{_size_f_i}, \code{_stride_f_i}\\
  spatial loop counters & \code{ctr_0}, \code{ctr_1}, \code{ctr_2}\\
  counter of sparse kernels & \code{sparse_idx}\\
  extracted constants & \code{__c_<value>} (e.g.\ \code{__c_1_0o4_0} for $1.0/4.0$)\\
  vector main loop / remainder loop & \code{ctr_i__1} / \code{ctr_i__0}, split at \code{ctr_i__rem_start}\\
  reduction accumulator & \code{<symbol>_local}\\
  \bottomrule
\end{tabular}
\end{center}

\begin{figure}[htbp]
  \centering
  % How field accesses become linear memory offsets (C layout, 5 x 4 array).
\begin{tikzpicture}[font=\footnotesize]
  % ---- the array (x = ctr_0 to the right, y = ctr_1 upwards) --------------------
  \begin{scope}[x=0.62cm,y=0.62cm]
    \foreach \x in {0,...,4} \foreach \y in {0,...,3} {
      \pgfmathtruncatemacro{\off}{4*\x+\y}
      \node[font=\tiny,text=psgray] at (\x+0.5,\y+0.5) {\off};
    }
    \cellfill[centercell]{2}{1} \cellfill[nbcell]{3}{1} \cellfill[nbcell]{2}{2}
    \node[font=\tiny,text=white] at (2.5,1.5) {9};
    \node[font=\tiny] at (3.5,1.5) {13};
    \node[font=\tiny] at (2.5,2.5) {10};
    \draw[gridline] (0,0) grid (5,4);
    \draw[axisarrow] (0,-0.4) -- (2,-0.4) node[right,celltext] {$x$ = \texttt{ctr\_0}};
    \draw[axisarrow] (-0.4,0) -- (-0.4,2) node[left,celltext] {$y$ = \texttt{ctr\_1}};
    \node[above,align=center,celltext] at (2.5,4.05) {\texttt{src}, shape $(5, 4)$, C layout\\
      strides: \texttt{\_stride\_src\_0} = 4, \texttt{\_stride\_src\_1} = 1};
  \end{scope}

  % ---- the formulas ------------------------------------------------------------
  \node[anchor=west,align=left,font=\scriptsize] at (4.2,1.6)
    {cell \texttt{(ctr\_0, ctr\_1) = (2, 1)}:\\[3pt]
     \texttt{src.center}\ \ $\to$ \texttt{\_data\_src[2*4 + 1*1]}\ = \textbf{9}\\
     \texttt{src[1,0]} (E) $\to$ \texttt{\_data\_src[(2+1)*4 + 1*1]} = \textbf{13}\\
     \texttt{src[0,1]} (N) $\to$ \texttt{\_data\_src[2*4 + (1+1)*1]} = \textbf{10}};

  % ---- memory --------------------------------------------------------------------
  \begin{scope}[shift={(0,-1.25)},x=0.6cm,y=0.5cm]
    \foreach \i in {0,...,19} {
      \pgfmathtruncatemacro{\col}{\i/4}
      \pgfmathsetmacro{\xx}{\i + 0.25*\col}
      \ifnum\i=9 \fill[centercell] (\xx,0) rectangle ++(1,1);
      \else\ifnum\i=13 \fill[nbcell] (\xx,0) rectangle ++(1,1);
      \else\ifnum\i=10 \fill[nbcell] (\xx,0) rectangle ++(1,1);
      \fi\fi\fi
      \draw[gridline] (\xx,0) rectangle ++(1,1);
      \node[font=\tiny] at (\xx+0.5,0.5) {\i};
    }
    \foreach \c in {0,...,4} {
      \pgfmathsetmacro{\xs}{4*\c + 0.25*\c}
      \draw[decorate,decoration={brace,mirror,amplitude=3pt}] (\xs,-0.1) -- ++(4,0)
        node[midway,below=3pt,font=\tiny] {$x=\c$};
    }
    \node[anchor=east,celltext] at (-0.2,0.5) {memory};
  \end{scope}
\end{tikzpicture}

  \caption[Linearized memory accesses]{Linearized memory accesses for a C-ordered $5\times4$ array: every access becomes \texttt{\_data\_src[ctr\_0 * \_stride\_src\_0 + ctr\_1 * \_stride\_src\_1]} plus the offsets. The north neighbour is the next element in memory, the east neighbour is one stride (4 elements) away.}
  \label{r01:fig:linearization}
\end{figure}

Further properties of the code:
\begin{itemize}
  \item Fields that are only read get \code{const} element types. All field pointers are
        \code{RESTRICT} (the JIT defines this macro as \code{__restrict__}), which
        promises the C compiler that fields do not alias.
  \item Memory accesses are \emph{linearized} (\cref{r01:fig:linearization}): \code{f[1,0]} becomes
        \code{_data_f[(ctr_0 + 1) * _stride_f_0 + ctr_1 * _stride_f_1]}. Fixed sizes
        and strides become literals.
  \item Every variable is declared \code{const} unless it is updated later, and every
        integer literal carries its type (\code{1LL} for \code{int64}).
  \item \code{get_c_code()} prints the function only. The \code{#include} directives for
        \code{required_headers} and the \code{RESTRICT} definition are added by whoever
        compiles it: the JIT, or a C++ code generator.
\end{itemize}

\paragraph{The IR view.}
\code{kernel.get_ir_code()} (or \code{ps.inspect(kernel)}) prints the same kernel in
the backend's typed intermediate representation. Every constant shows its type, and
every declaration shows the symbol's type. This view is most useful when you debug type
errors:
\irsnippet[nolol=false,caption={The Jacobi kernel as typed IR (\texttt{kernel.get\_ir\_code()}).}]{ex01_first_kernel.ir}

\section{Just-in-time compilation}\label{r01:sec:jit}

\code{kernel.compile()} delegates to the kernel's JIT compiler, an instance of
\cls{JitBase} (\file{src/pystencils/jit/}). The JIT comes from \code{config.jit}. If
that option is unset, it is chosen by target: \cls{CpuJit} for CPU targets, \cls{CupyJit}
for CUDA and HIP (or \code{no_jit} if CuPy is not installed), and \cls{SYCLJit} for SYCL.
\code{compile()} returns a \cls{KernelWrapper}. That is a callable that takes keyword
arguments and also exposes \code{.kernel}, \code{.parameters} and \code{.code}.

\subsection{How the CPU JIT works}
\cls{CpuJit} combines two components: an \term{extension module builder}, which renders
C++ source code, and a \cls{CompilerInfo}, which describes the host compiler.
\begin{enumerate}
  \item \textbf{Render.} \cls{DefaultExtensionModuleBuilder} fills the template
        \file{jit/cpu/cpu_kernel_module.tmpl.cpp}. The result is a complete CPython
        extension module of about 290 lines: the kernel in \code{namespace internal},
        NumPy array proxies, argument checks and a module function \code{invoke}.
  \item \textbf{Hash.} The module name is \code{module_<sha256>}. The hash covers the
        source code and the full compiler command line.
  \item \textbf{Compile.} If no matching shared object is in the cache, it runs
        \code{<cxx> <cxxflags> -I<python> -I<pystencils/include> -I<numpy> -shared -o
        module.so module.cpp}. A file lock serializes concurrent compilations.
  \item \textbf{Load.} The shared object is imported with \code{importlib}. The
        returned wrapper forwards \code{**kwargs} to \code{module.invoke} (class
        \cls{DefaultCpuKernelWrapper}).
\end{enumerate}
The interesting part of the rendered module is the argument handling.
\Cref{r01:lst:glue} shows it for the kernel \code{g = 2 f}. Each parameter is filled
according to its property: \code{FieldBasePtr} from the array's data pointer,
\code{FieldShape} and \code{FieldStride} from the array's shape and strides (converted
from bytes to elements), and scalars from Python numbers.
\cppsnippet[nolol=false,caption={Argument extraction in the JIT extension module (excerpt).},label=r01:lst:glue,basicstyle=\ttfamily\scriptsize]{ex10_jit.module-glue}

\subsection{Calling a compiled kernel}
Arguments are passed \emph{by keyword only}: field names for arrays, symbol names for
scalars. Before the kernel runs, the wrapper checks the arguments and reports problems
as Python exceptions:
\outsnippet[basicstyle=\ttfamily\scriptsize]{ex10_jit.call-errors}
There is no implicit conversion: a \code{float32} array passed to a \code{double} field
is an error. Arrays are passed \emph{including} their ghost layers. The kernel decides
from its iteration space which cells to update. All \code{GENERIC} domain fields must
have the same spatial shape.

\subsection{Choosing the compiler and flags}
The default \cls{CompilerInfo} is \cls{GccInfo} on Linux (if \texttt{g++} works),
otherwise \cls{ClangInfo}. On macOS it is \cls{AppleClangInfo}. Its fields are
\code{openmp=True}, \code{optlevel="fast"}, \code{cxx_standard="c++17"},
\code{target=Target.CurrentCPU} and \code{extra_cxxflags=[]}. A custom JIT is created
and passed like this:
\pysnippet{ex10_jit.custom-jit}
\outsnippet[basicstyle=\ttfamily\scriptsize]{ex10_jit.flags}
The \code{-march} flag comes from the \cls{CompilerInfo}'s \emph{own} target. For
\code{CurrentCPU} that is \code{-march=native}, except on RISC-V, where GCC rejects
\code{native} and \ps{} synthesizes \code{-march=rv64gcv_zvl256b} or
\code{-march=rv64gc}. Explicit targets map to \code{x86-64-v2/v3/v4},
\code{armv8-a}, \code{armv8.2-a+fp16}, \code{armv8.2-a+sve} and
\code{rv64gcv_zvl256b}. \cls{ClangInfo} replaces \code{-Ofast} by
\code{-O3 -ffast-math -fhonor-infinities}. The compiler executables \texttt{g++} and
\texttt{clang++} are hard-coded in \code{cxx()}. To use another binary (such as a
specific GCC version), subclass the info class and override \code{cxx()}.

\begin{pitfall}[The JIT does not follow the kernel's target]
  \code{CompilerInfo.target} is independent of \code{CreateKernelConfig.target}: the
  default JIT always compiles with \code{-march=native}. A vectorized
  \code{X86_AVX512} kernel compiled on this AVX2 laptop therefore fails with
  \emph{``inlining failed in call to `always\_inline' `\_mm512\_storeu\_pd': target
  specific option mismatch''}. The kernel can only run on matching hardware anyway. To
  compile for a target other than the host (e.g.\ to inspect the assembly), construct
  the \cls{CompilerInfo} with that target, as shown for RVV in \cref{r01:ch:rvv-backend}.
\end{pitfall}

\paragraph{Cache.}
\code{CpuJit(objcache=AUTO)} (the default) caches modules in
\file{~/.cache/pystencils/cpujit} (via \code{appdirs}). \code{objcache=None} compiles in
a temporary directory, and a path selects a custom cache. \code{emit_warnings=True}
turns compiler warnings into Python \code{RuntimeWarning}s.

\section{GPU execution with CuPy}\label{r01:sec:gpu-execution}

\cls{CupyJit} (\file{jit/gpu_cupy.py}) compiles a \code{GpuKernel} at run time with
NVRTC through \code{cupy.RawKernel}. It prepends the required headers and
\code{#define RESTRICT __restrict__}, and compiles with \code{-std=c++17} plus the \ps{}
include path. The wrapper \cls{CupyKernelWrapper} accepts CuPy arrays as keyword
arguments and checks their types, shapes and strides. It then evaluates the kernel's
launch configuration (\cref{r01:sec:gpu-options}) on the actual array sizes to obtain block
and grid, and launches the kernel on the arrays' device. The launch configuration is
available as \code{kfunc.launch_config}. That is the place to call
\code{fit_block_size(...)}, or to set \code{block_size}/\code{grid_size} for manual
launch grids. The GPU examples in this manual were generated without a GPU: code
generation does not need CuPy.

\section{Using kernels outside Python}\label{r01:sec:external-use}

A \code{Kernel} object is a complete, self-describing unit, so other tools can generate
C++ code around it. Such tools set \code{config.jit = ps.no_jit} (there is no need to
compile in Python) and then use three things from the kernel:
\begin{enumerate}
  \item \code{kernel.get_c_code()} for the function definition, and
        \code{kernel.required_headers} for the includes, plus the \ps{} runtime headers
        found with \code{python -m pystencils.include};
  \item \code{kernel.parameters} with their properties, to map C++ data structures onto
        the arguments;
  \item \code{kernel.metadata} (e.g.\ ghost layers) and, for GPU kernels, the launch
        configuration.
\end{enumerate}

\paragraph{pystencils-sfg and waLBerla.}
waLBerla's new code generation layer \emph{sweepgen} (in
\file{submodules/walberla/sweepgen}) is built on pystencils-sfg:
\begin{itemize}
  \item A CMake function \code{walberla_generate_sources(<target> SCRIPTS gen.py)} runs
        a Python generator script at build time.
  \item Inside the script, \code{with SourceFileGenerator() as sfg:
        sfg.generate(Sweep(name, assignments, config))} describes a sweep.
        \code{Sweep.generate} creates the \ps{} kernel through
        \code{sfg.kernel_namespace(...).create(assignments, name, config)}. The
        configuration comes from \code{get_pystencils_config()}
        (\file{sweepgen/build_config.py}), which sets the data type, the target and
        OpenMP, and \code{jit = no_jit}.
  \item For every field parameter, \code{sfg.map_field(field, <extraction>)} generates
        C++ code that takes the data pointer, shapes and strides from a waLBerla
        \code{GhostLayerField}. It relies on exactly the \code{FieldBasePtr},
        \code{FieldShape} and \code{FieldStride} properties shown in \cref{r01:lst:params}.
  \item The result is a C++ class with a properties struct and an \code{operator()} that
        waLBerla calls on every block of the distributed domain.
\end{itemize}
The older \file{python/pystencils_walberla} path in waLBerla targets the 1.x API (it
imports \code{pystencils.backends.simd_instruction_sets}, which no longer exists) and
does not work with \ps{}~2.0.

\section{Runtime helpers for Python simulations}

For prototyping whole simulations in Python, \ps{} ships some convenience layers on top
of the JIT:
\begin{description}
  \item[\code{ps.create_data_handling(domain_size, periodicity, default_ghost_layers, ...)}]
        returns a \cls{SerialDataHandling} that allocates arrays with ghost layers
        (\code{add_array}, returning a matching fixed-size field), fills them, swaps
        them, synchronizes ghost layers periodically, runs kernels with
        \code{dh.run_kernel(kfunc)} (passing all arrays by name), and moves data
        between CPU and GPU. A \code{ParallelDataHandling} for distributed runs needs
        waLBerla's Python bindings.
  \item[\code{pystencils.boundaries}] provides \cls{BoundaryHandling} with
        \cls{Neumann} and \cls{Dirichlet} conditions. Internally, these are sparse index
        kernels (\cref{r01:sec:sparse}).
  \item[\code{pystencils.runhelper}] offers a \cls{ParameterStudy} driver and a
        result \cls{Database}.
\end{description}

\begin{keyidea}[Summary]
  \begin{itemize}[leftmargin=1.2em]
    \item A \code{Kernel} is one C function plus metadata: parameters with properties
          (\code{FieldBasePtr}, \code{FieldShape}, \code{FieldStride}), headers, ghost
          layers.
    \item \code{kernel.compile()} builds and loads a Python extension module; call it with
          keyword arguments, passing whole arrays including ghost layers.
    \item The JIT compiles for the host (\code{-march=native}) unless you give it a
          \code{CompilerInfo} with another target.
    \item For C++ applications, generate with \code{no_jit} and let pystencils-sfg or
          waLBerla map their data structures onto the parameters.
  \end{itemize}
\end{keyidea}
